\textbf{Corollary 5 (multiplicity two).} For \(K\geq11\), form all signed unordered pairs with distinct time labels: \[
 \epsilon v_t+\delta w_s,
 \quad 0\leq t<s<K,\quad v,w\in\{u,z\},\quad
 \epsilon,\delta\in\{-1,1\}.
\] There are \(8K(K-1)\) such pairs. Their vector keys are nonzero, have multiplicity at most two, and have precisely \[
 4K\ \hbox{double keys},\quad 8K(K-2)\ \hbox{single keys},
 \quad 8K^2-12K\ \hbox{distinct keys}.                          \tag{CB.9}
\] The double classes are the two identities \[
 z_{j+1}-z_j=u_{j+1}+u_j,\qquad
 z_{j+1}-u_j=u_{j+1}+z_j,
\] with their negatives, for j modulo K. The third partition of (CB.8) uses equal-time pairs and is not in the domain. Each displayed signed pair belongs to only one class: its species, the two adjacent cyclic labels, and its signs determine the other pair. Thus there are no triple classes or overlapping equivalence chains. Lemma 4 ensures that no additional modular relation is missed. K=3 is a useful negative control: it has six triple keys, and the range restriction must not be removed.

\subsection{Weighted phase control and a finite normal form for C}\label{coboundary:weighted-phase-control-and-a-finite-normal-form-for-c}

Write \(C_j=\widetilde U_0^jC\widetilde U_0^{-j}\), \(\bar C=K^{-1}\sum_j C_j\), \(G_C=\sum_j c_jC_j\), where \(c_j=(K-1-2j)/(2K)\). Then \[
 G_C-\widetilde U_0G_C\widetilde U_0^*=C-\bar C,\qquad [\bar C,\widetilde U_0]=0.                \tag{CB.10}
\] Define a phase majorant, including both species, \[
 a_j=\max_{v,w\in\{u,z\}}
          |e_N(\omega(v,T^jw))-1|.
\] One has \(a_{K-j}=a_j\). In particular, no independence assumption about the several phase factors is being made.

For a centered representative \(\delta\) of j modulo K, \[
 \omega(v,T^\delta w)
 =(2v_1w_1-v_2w_2)P_{2\delta}+\det(v,w)Q_{2\delta}.
\] For the four choices of species these are \(-8P_{2\delta}\), \(16P_{2\delta}\), and \(\pm8Q_{2\delta}\), respectively. Their absolute values are at most \(8\lambda^{|\delta|}\). If \(|\delta|=(K-1)/2-\ell>0\), this gives \[
 a_j\leq\frac{32\pi}{\alpha(1-\lambda^{-K})}\lambda^{-\ell}.
\] At zero, \(a_0\leq16\pi/N\). Summing both signs, including their weights \(\ell/K\) and \((\ell+1)/K\), proves the explicit estimates \[
 \begin{split}
 L_K:=\sum_j a_j
   &\leq\frac{32\pi}{1-\lambda^{-K}}+\frac{16\pi}{N}<128,\\
 K\sum_j|c_j|a_j
   &\leq\frac{16\sqrt2\pi}{1-\lambda^{-K}}+
                                      \frac{8\pi K}{N}<128.
 \end{split}                                                    \tag{CB.11}
\] For \(K\geq3\), use \(\lambda^{-K}<1/10\), \(N\geq2K+1\), \(\pi<22/7\), and \(\sqrt2<3/2\) to check the last inequalities. These constants are deliberately conservative.

The commutator \([C_t,C_s]\) has sixteen Weyl terms, each of modulus at most \(a_{s-t}\). Corollary 5 says that each key in the sum over \(t<s\) has at most two preimages. Weyl matrices are orthonormal in the \emph{full} normalized Hilbert--Schmidt inner product, also over composite odd N. Cauchy--Schwarz inside each key class therefore gives \[
 \begin{split}
 \|[G_C,\bar C]\|_{2,N}^2
 &\leq\frac{32}{K^4}\sum_{j=1}^{K-1}(K-j)j^2a_j^2\\
 &=\frac{16}{K^3}\sum_{j=1}^{K-1}j(K-j)a_j^2
 \leq\frac{1024}{K}.
 \end{split}                                                    \tag{CB.12}
\] The equality between the two sums uses \(a_{K-j}=a_j\); the final bound uses \(a_j^2\leq2a_j\) and (CB.11). Norm restriction to the odd subspace obeys \(\|M\|_{2,-}\leq\sqrt{2N/(N-1)}\|M\|_{2,N}\) by positivity. We have not assumed orthogonality of restricted Weyl matrices.

The triangle inequality in operator norm also gives \[
 \|[G_C,C]\|_{\rm op}\leq2048/K,\qquad
 \|[C,\bar C]\|_{\rm op}\leq2048/K.                           \tag{CB.13}
\] The usual two-exponential Duhamel inequality is \(\|e^{ixA}e^{ixD}-e^{ix(A+D)}\|\leq x^2\|[A,D]\|/2\) in any unitarily invariant norm. Applying it first to -\(G_{C}\),C, then to C-\(G_{C}\) and \(\widetilde U_0G_C\widetilde U_0^*\), using (CB.10)--(CB.13), proves \[
 \left\|e^{-i\zeta G_C}(e^{i\zeta C}\widetilde U_0)e^{i\zeta G_C}
                        -e^{i\zeta\bar C}\widetilde U_0\right\|_{2,-}
 \leq\frac{1000\zeta^2}{\sqrt K}\quad(K\geq11).               \tag{CB.14}
\] Indeed the less rounded bound is \(\zeta^2(2048/K+16\sqrt{2N/(N-1)}/\sqrt K)\), which is smaller than (CB.14) already using \(K\geq11\) and \(N\geq7\). This is another estimate valid for every real zeta, not a formal expansion requiring a small generator.

Combining the two gauges and telescoping unitary powers gives \[
 \left|\tau_-(V_\varepsilon^{mK})-
          \tau_-(e^{im\zeta K\bar C})\right|
 \leq |m|KR_N(\varepsilon)+1000|m|\zeta^2\sqrt K.              \tag{CB.15}
\]

\subsection{The two-species Gaussian limit}\label{coboundary:the-two-species-gaussian-limit}

\textbf{Lemma 6.} Put \(X_K=(\sqrt K/2)\bar C\). For every fixed \(r\geq1\), \[
 \tau_-(X_K^{2r})=(2r-1)!!+O_r(K^{-1})+O_r(K^r/N).             \tag{CB.16}
\] For every fixed odd m, eventually \(|\tau_-(X_K^m)|\leq(4K)^{m/2}/(N-1)\). Consequently its empirical spectral measures converge weakly and in quadratic Wasserstein distance to standard normal.

\textbf{Proof.} There are 4K signed Weyl atoms, with coefficient +1 for species u and -1 for species z. For a word with central phase \(e_{N}\)(\(\phi\)) and total index v, the exact odd trace is \[
 e_N(\phi)\frac{N\mathbf1_{v=0}-1}{N-1}.                       \tag{CB.17}
\] This follows from \(\operatorname{Tr}W(v)=N\mathbf1_{v=0}\) and \(\operatorname{Tr}(RW(v))=1\). The noncentral contribution after normalization \((4K)^{-m/2}\) is at most \((4K)^{m/2}/(N-1)\).

Apply Lemma 4 to central words. The following bounded-cluster argument handles the unavoidable identities (CB.8), rather than incorrectly excluding them. In a nonzero polynomial relation with at most m signed atoms and coefficients 1 or \(\sqrt2\), a gap d with \(\lambda^d>2m^2\) separates independently vanishing components. To see this, write \(P=P_{\rm low}+Z^{a+d}P_{\rm high}\) with degree of \(P_{low}\) at most a. If \(P_{\rm high}(\lambda)\ne0\), its nonzero integer norm and the bound \(|P_{\rm high}(\lambda^{-1})^{\rm conjugate}|\leq\sqrt2m\) give \(|P_{\rm high}(\lambda)|\geq1/(\sqrt2m)\). This contradicts cancellation by \(|P_{\rm low}(\lambda)|\leq\sqrt2m\lambda^a\). Iterating produces components with spans bounded solely in m.

Cancel opposite identical atoms first and restore these pairs after splitting. Every nonzero component then contains at least three occurrences. A two-atom relation cannot change species, because its two algebraic norms would be 1 and -2, and cannot change exponent, because \(\lambda>1\). A nonbalanced 2r-word thus has at most r-1 components. There are only finitely many partitions of its positions, species, signs and bounded offset patterns at fixed r. Each component has K possible cyclic translations, so there are \(O_r(K^{r-1})\) such words. Choosing the largest-gap cut introduces no additional K factor: this count is over the actual translated components. Their normalized contribution is \(O_r(1/K)\).

In a balanced word, pair identical species-time labels with opposite signs. For each pairing there are \(K^r\) time choices, \(2^r\) species choices and \(2^r\) sign choices. Their product is precisely the normalization \((4K)^r\). Repeated pair labels cause at most \(O_r(K^{r-1})\) overcount. Signs from species z occur twice in each pair and hence multiply to +1. The Weyl phase is a product of crossing factors or their conjugates. For fixed species, each crossing time difference is uniform on the cyclic group. By (CB.11) and \(|\prod u_j-1|\leq\sum|u_j-1|\), its average error is at most 128 times the number of crossings divided by K. No independence of the crossing factors is assumed. Summing over the \((2r-1)!!\) pairings gives (CB.16).

For odd m, reduce a lifted relation in \(\mathbb Z[\sqrt2]\) modulo 2. Here \(\lambda=3+2\sqrt2\) reduces to 1. The coefficients of both 1 and \(\sqrt2\) in the relation must be even. Both species counts would therefore be even, contradicting odd m. Thus no odd central word exists after lifting.

Bounded second moments yield tightness, higher even moments yield uniform integrability, and the Gaussian moment sequence is determinate, since its Carleman sum diverges. Therefore all weak subsequential limits are normal. Convergence of the second moments gives the claimed Wasserstein upgrade. End of proof.

Since \(\zeta K\bar C=t_KX_K\), (CB.15) and Lemma 6 prove Theorem 1. More explicitly, the finite comparison error is bounded by \[
